\section{Introduction}
A reasoning model may verify a derivation, compare explanations, or discard an unoriginal research idea. These activities share operations of checking, exploration, and revision, but organize them for different purposes. This question of organization has a long history: \citet{minsky1986society} describes intelligence as arising from interacting simpler processes, while human neuroscience shows shared control networks reconfiguring across task demands \citep{cole2013hubs}. For language models, the corresponding empirical question is how observable reasoning operations vary across tasks, and what happens when their organization is changed.

Recent work characterizes reasoning behaviors, temporal structure, and model-specific patterns \citep{gandhi2025habits,chen2025lot,li2025thinkarm,lee2026reasonops}, connects conversational organization to learning \citep{kim2026society}, and modifies reasoning through steering and reward-trained vectors \citep{chen2025seal,sinii2025vectors}. Together, these lines leave a task-dependent question: when does an organization characteristic of one domain help another, and is a naturally occurring profile a useful intervention target?

We investigate this question with \emph{Thought Atlas}. A \emph{reasoning signature} is the prevalence and temporal allocation of annotated operations in emitted reasoning text. It describes observable behavior without assuming exhaustive access to internal computation. The atlas spans math, coding, science, planning, moral advice, idea generation, harmful requests, and cybersecurity. Within each model, we test whether operation amounts identify domains beyond trace length and whether timing supplies additional information.

The differences concern both allocation and function. Mathematical traces combine early decomposition with increasing verification; idea-generation traces emphasize alternatives and candidate revision. Planning examples check action preconditions, whereas moral examples use subgoals to organize advice. Thus, the same annotation can describe different work across domains. Matched instructions further test whether profiles change with what a model is asked to do with the same problem, alongside a separate output-format manipulation.

On Qwen3.5-9B, a common behavior-to-activation map constructs interventions from five domain profiles. The source--target matrix reveals cross-domain benefits, substantial asymmetry, and budget dependence. A selected within-math archetype provides a separate target, which we compare with a residual bias trained using a correctness-within-budget reward. Their behavioral similarity and performance differences distinguish the measured profile from its implementation and consequences. Idea generation then tests how interventions associated with objective-task benefits affect originality and feasibility.

Our contribution connects three findings: reproducible domain differences in operation amount and timing across seven models; model-dependent responses to matched task demands; and task-, budget-, and objective-dependent consequences of steering and reward learning. The atlas provides a common behavioral description across these settings. It establishes neither a universally beneficial signature nor causal mediation by a particular annotation.
